\documentclass[11pt,a4paper]{article}

% ================ Encodage et langue ================
\usepackage[utf8]{inputenc}
\usepackage[T1]{fontenc}
\usepackage[french]{babel}
\usepackage{lmodern}
\usepackage{textcomp}

% ================ Mise en page ================
\usepackage[margin=2.2cm]{geometry}
\usepackage{parskip}
\setlength{\parindent}{0pt}
\setlength{\parskip}{0.6em}

% ================ Couleurs et liens ================
\usepackage[table]{xcolor}
\definecolor{primary}{HTML}{0B5394}
\definecolor{accent}{HTML}{1F77B4}
\definecolor{light}{HTML}{F2F2F2}
\definecolor{codebg}{HTML}{F6F8FA}
\definecolor{rulec}{HTML}{D0D7DE}

\usepackage[colorlinks=true, linkcolor=primary, urlcolor=primary, citecolor=primary]{hyperref}

% ================ Tableaux ================
\usepackage{longtable, booktabs, array, tabularx}
\renewcommand{\arraystretch}{1.25}
\arrayrulecolor{rulec}

% ================ Listes ================
\usepackage{enumitem}
\setlist{itemsep=2pt, topsep=4pt}

% ================ Code et citations ================
\usepackage{fancyvrb}
\usepackage{framed}
\definecolor{shadecolor}{named}{light}

% Encadrés pour les blocs > (quotes)
\usepackage{mdframed}
\newmdenv[
  linecolor=primary,
  linewidth=2pt,
  topline=false, bottomline=false, rightline=false,
  innerleftmargin=10pt, innerrightmargin=8pt,
  innertopmargin=6pt, innerbottommargin=6pt,
  backgroundcolor=codebg,
  skipabove=8pt, skipbelow=8pt
]{quotebox}
\renewenvironment{quote}{\begin{quotebox}}{\end{quotebox}}

% ================ Titres ================
\usepackage{titlesec}
\titleformat{\section}{\Large\bfseries\color{primary}}{\thesection.}{0.6em}{}
\titleformat{\subsection}{\large\bfseries\color{accent}}{\thesubsection}{0.6em}{}
\titleformat{\subsubsection}{\normalsize\bfseries}{\thesubsubsection}{0.6em}{}

% ================ En-têtes et pieds de page ================
\usepackage{fancyhdr}
\pagestyle{fancy}
\fancyhf{}
\fancyhead[L]{\small\color{gray}ESPRIT School of Business --- M1 BA --- Projet Intégré}
\fancyhead[R]{\small\color{gray}\nouppercase{\leftmark}}
\fancyfoot[C]{\small\thepage}
\renewcommand{\headrulewidth}{0.3pt}
\renewcommand{\footrulewidth}{0pt}

% ================ Pandoc helpers ================
\providecommand{\tightlist}{\setlength{\itemsep}{0pt}\setlength{\parskip}{0pt}}
\providecommand{\passthrough}[1]{#1}

\usepackage{calc}
\providecommand{\real}[1]{#1}


% ================ Métadonnées ================
\title{\color{primary}\textbf{Projet Intégré --- Analyse et Prédiction
du Churn Client}}
\author{Aymen Ben Brik}
\date{2026}

\begin{document}

\begin{titlepage}
\centering
\vspace*{1.5cm}

{\Large\bfseries ESPRIT School of Business}\\[0.3cm]
{\large Master 1 --- Business Analytics (M1 BA)}\\[2cm]

{\color{gray}\large Projet Intégré}\\[0.3cm]
{\color{gray}\large Analyse et Prédiction du Churn Client}\\[2.2cm]

\rule{0.85\linewidth}{1.2pt}\\[0.6cm]
{\Huge\bfseries\color{primary} Projet Intégré --- Analyse et Prédiction
du Churn Client\par}
\rule{0.85\linewidth}{1.2pt}\\[2cm]

{\large Tuteur du projet}\\[0.2cm]
{\Large\bfseries Aymen Ben Brik}\\[1cm]
{\large 2026}

\vfill
{\color{gray}\small Document à destination des étudiants}
\end{titlepage}

\tableofcontents
\newpage

\section{Projet Intégré --- Analyse et Prédiction du Churn
Client}\label{projet-intuxe9gruxe9-analyse-et-pruxe9diction-du-churn-client}

\begin{quote}
\textbf{École} : ESPRIT School of Business \textbf{Programme} : Master 1
--- Business Analytics (M1 BA) \textbf{Tuteur du projet} : Aymen Ben
Brik \textbf{Année} : 2026
\end{quote}

\subsection{1. Contexte}\label{contexte}

Une institution bancaire met à votre disposition un jeu de données réel
et anonymisé décrivant ses clients, leurs comptes et leurs produits. La
fidélisation client (lutte contre l'attrition, ou \emph{churn}) est un
enjeu majeur du secteur : acquérir un nouveau client coûte plusieurs
fois plus cher que d'en conserver un existant.

Votre mission est de concevoir, de bout en bout, une \textbf{solution
analytique complète} permettant de comprendre, mesurer et anticiper le
churn, puis de mettre ces analyses à disposition des décideurs métier
via une interface web.

\subsection{2. Objectifs pédagogiques}\label{objectifs-puxe9dagogiques}

Ce projet vous permet de mettre en pratique, sur un cas réel, l'ensemble
de la chaîne décisionnelle :

\begin{itemize}
\tightlist
\item
  Concevoir une architecture de données analytique (modélisation
  dimensionnelle, ETL).
\item
  Exploiter les techniques de Business Intelligence pour produire des
  indicateurs métier.
\item
  Construire et évaluer un modèle de Machine Learning supervisé.
\item
  Restituer les résultats via des tableaux de bord et une application
  web déployée.
\item
  Travailler en équipe avec une démarche projet structurée (versioning,
  livrables, présentation).
\end{itemize}

\subsection{3. Périmètre fonctionnel}\label{puxe9rimuxe8tre-fonctionnel}

Le projet est structuré en \textbf{cinq composantes complémentaires} qui
constituent ensemble une chaîne décisionnelle complète :

\subsubsection{3.1 ETL (Extract -- Transform --
Load)}\label{etl-extract-transform-load}

\begin{itemize}
\tightlist
\item
  Extraire les données brutes (fichiers CSV et tables de dimensions
  Excel).
\item
  Nettoyer, déduplicaster, gérer les valeurs manquantes et les
  incohérences.
\item
  Transformer et enrichir les données (jointures avec les dimensions,
  calculs dérivés).
\item
  Charger les données préparées dans un entrepôt analytique (Data
  Warehouse).
\end{itemize}

\subsubsection{3.2 Business Intelligence
(BI)}\label{business-intelligence-bi}

\begin{itemize}
\tightlist
\item
  Concevoir un modèle dimensionnel en étoile (faits + dimensions).
\item
  Définir les KPIs métier (taux de churn, ancienneté moyenne, solde
  moyen, segmentation, etc.).
\item
  Mettre en place les requêtes analytiques nécessaires aux dashboards.
\end{itemize}

\subsubsection{3.3 Machine Learning}\label{machine-learning}

\begin{itemize}
\tightlist
\item
  Préparer un jeu de données d'entraînement à partir du Data Warehouse.
\item
  Entraîner et comparer plusieurs modèles de classification supervisée
  pour prédire la probabilité de churn d'un client.
\item
  Évaluer les modèles avec des métriques adaptées à un problème
  déséquilibré.
\item
  Interpréter les variables explicatives les plus importantes.
\item
  Sauvegarder le modèle final pour intégration.
\end{itemize}

\subsubsection{3.4 Power BI}\label{power-bi}

\begin{itemize}
\tightlist
\item
  Produire un ou plusieurs rapports Power BI alimentés par le Data
  Warehouse.
\item
  Restituer les KPIs et permettre une exploration interactive (filtres,
  drill-down).
\item
  Au minimum : vue d'ensemble, analyse du churn, segmentation client.
\end{itemize}

\subsubsection{3.5 Interface Web de
déploiement}\label{interface-web-de-duxe9ploiement}

\begin{itemize}
\tightlist
\item
  Développer une application web qui expose le modèle ML et les analyses
  BI.
\item
  L'utilisateur doit pouvoir :

  \begin{itemize}
  \tightlist
  \item
    Saisir les caractéristiques d'un client et obtenir une prédiction de
    churn.
  \item
    Visualiser la liste des clients actifs à risque élevé.
  \item
    Consulter des indicateurs de synthèse.
  \end{itemize}
\item
  Déployer l'application sur une plateforme cloud accessible
  publiquement.
\end{itemize}

\subsection{4. Livrables attendus}\label{livrables-attendus}

Trois livrables sont attendus à la fin du mois :

{\def\LTcaptype{none} % do not increment counter
\begin{longtable}[]{@{}
  >{\raggedright\arraybackslash}p{(\linewidth - 4\tabcolsep) * \real{0.3333}}
  >{\raggedright\arraybackslash}p{(\linewidth - 4\tabcolsep) * \real{0.3333}}
  >{\raggedright\arraybackslash}p{(\linewidth - 4\tabcolsep) * \real{0.3333}}@{}}
\toprule\noalign{}
\begin{minipage}[b]{\linewidth}\raggedright
Livrable
\end{minipage} & \begin{minipage}[b]{\linewidth}\raggedright
Format
\end{minipage} & \begin{minipage}[b]{\linewidth}\raggedright
Contenu attendu
\end{minipage} \\
\midrule\noalign{}
\endhead
\bottomrule\noalign{}
\endlastfoot
\textbf{Code source} & Dépôt GitHub public & Notebooks, scripts ETL,
code ML, application web, README détaillé, instructions d'installation
et d'exécution \\
\textbf{Présentation orale} & Slides (PPT/PDF), 20 min + 10 min de
questions & Démonstration de l'application, présentation de la démarche,
principaux résultats \\
\textbf{Rapport} & PDF, 25--40 pages & Document professionnel décrivant
la démarche, l'architecture, les choix techniques, les résultats, les
limites et les perspectives \\
\end{longtable}
}

\subsection{5. Modalités}\label{modalituxe9s}

\begin{itemize}
\tightlist
\item
  \textbf{Composition des équipes} : 3 à 4 étudiants par groupe.
\item
  \textbf{Durée} : 1 mois (4 semaines), à temps partagé avec les autres
  cours.
\item
  \textbf{Encadrement} : un point d'avancement hebdomadaire avec
  l'encadrant.
\item
  \textbf{Outils} : libre choix dans le respect des contraintes énoncées
  dans le guide étudiant (document 4).
\end{itemize}

\subsection{6. Critères d'évaluation}\label{crituxe8res-duxe9valuation}

{\def\LTcaptype{none} % do not increment counter
\begin{longtable}[]{@{}ll@{}}
\toprule\noalign{}
Critère & Pondération \\
\midrule\noalign{}
\endhead
\bottomrule\noalign{}
\endlastfoot
Qualité de l'ETL et du modèle dimensionnel & 15 \% \\
Pertinence et qualité des analyses BI / Power BI & 20 \% \\
Rigueur de la démarche Machine Learning & 20 \% \\
Qualité de l'interface web et du déploiement & 15 \% \\
Qualité du code (lisibilité, structure, documentation) & 10 \% \\
Qualité du rapport écrit & 10 \% \\
Qualité de la présentation orale & 10 \% \\
\end{longtable}
}

\subsection{7. Règles importantes}\label{ruxe8gles-importantes}

\begin{itemize}
\tightlist
\item
  \textbf{Confidentialité} : les données sont confidentielles. Elles ne
  doivent \textbf{pas} être publiées sur GitHub (utilisez un
  \texttt{.gitignore}). Seul le code peut être partagé publiquement.
\item
  \textbf{Anonymat} : ne mentionnez ni le nom de l'institution dont
  proviennent les données, ni d'éléments permettant de l'identifier, ni
  dans le code, ni dans le rapport, ni dans la présentation.
\item
  \textbf{Versioning} : le dépôt GitHub doit montrer un historique
  régulier de commits par chaque membre de l'équipe.
\item
  \textbf{Reproductibilité} : tout doit pouvoir être ré-exécuté par un
  tiers à partir du README.
\end{itemize}

\end{document}
